% Final TCP, UDP and block measurements belong in the marked result slots.
% Earlier TCP observations below are scoped to CLAIMS S1, S3, F5/F6 and F10.
\section{Evaluation}
\label{sec:eval}

The central question is how much CPU a service spends delivering a byte that
is already in memory. TCP separates the cost of capturing, sending and
forwarding that byte (\S\ref{sec:eval:tcp}). UDP tests whether the result holds
when every datagram requires a separate routing decision
(\S\ref{sec:eval:udp}). Block I/O tests whether the same object can move between
storage and a socket efficiently in either direction (\S\ref{sec:eval:block}).
We also examine the cost of keeping objects when clients run at different speeds
(\S\ref{sec:eval:retention}). The measurements quoted below come from earlier
TCP prototypes. The marked figures reserve the corresponding comparisons for the
final builds, including UDP and block I/O.

\subsection{Method}
\label{sec:eval:method}

The earlier network testbed used cloud VMs with virtio-net. For relay workloads,
a source sends through the machine under test to a separate sink. The relay
uses different interfaces for the two legs. The primary metric is CPU-seconds on the host under test per gigabyte of
useful payload. For network delivery we count bytes at the receiver. For block
ingest we count bytes in completed writes. Goodput is
reported alongside CPU cost: a link-limited service can use less CPU without
delivering more bytes per second. Request/response workloads also report
latency and CPU per completed transaction.
\todom{Give final CPU and memory configuration, kernel builds, NIC and storage
devices, network and device ceilings, worker/interrupt placement, repeats and
uncertainty procedure for each testbed.}

A matched comparison uses the same application work, request size, batching,
outstanding depth and CPU placement. Setup and warm-up are outside the steady
measurement window on both sides. Host accounting includes interrupt, helper
and reclamation work. The final TCP bulk comparisons run at full speed, with
fixed-rate tests to isolate CPU cost at equal goodput. CPU is computed from host counters with a correction for
idle time that overlaps interrupt activity. Results from different estimators or
distinct throughput regimes are kept separate. A verification run on the same
path checks the delivered data outside the timed comparison. Where an interface
changes ownership or does different work, the text identifies that difference.

\subsection{TCP: capture, send and forward}
\label{sec:eval:tcp}
\label{sec:eval:rx}
\label{sec:eval:ubiquitous}
\label{sec:eval:tx}
\label{sec:eval:fwd}

\paragraph{Capture.} A copying receiver and a named receiver can both retain a
complete stream range. A mapped receiver can retain the same range as an object
in the I/O area. A raw mapped receiver instead acquires and returns receive
slots, and a discard receiver does neither the retention nor the subsequent
read. These last two paths show the cost of different services. The comparison
holds source writes fixed and records the bytes actually returned to the
receiver, since TCP can split or coalesce writes.

\paragraph{Send.} A prepared-data test removes capture from the timed path.
The baselines are copying sends and zero-copy sends with either a per-send pin
or a pre-registered pool. The registered sender keeps its source unchanged, so
it can continue sending while reuse notifications arrive. Comparing the two
\sendzc paths at equal batching and depth measures the effect of pinning once
at registration. A separate test sends captured objects
with the fragmentation that the receive path actually produces.

On an earlier prototype, full-speed named sends used 15--27\% less sending-host
CPU per delivered byte than matched \sendzc requests that pinned pages per send,
at the tested sizes from 4\KiB to 1\MiB. There were three runs per size and path.
A different comparison on that prototype used batches of eight ranges on both
sides and a registered, never-rewritten buffer. At 1 and 4\KiB under saturation,
named sends cost 6--13\% more CPU per delivered byte. Which cost dominates
depends on the work the baseline already amortizes.

\paragraph{Forward.} A relay receives a range and sends it on a second TCP
connection. The alternatives include copying through \iouring and \splice
through a pipe. All preserve byte order and backpressure. A mapped-object relay
lets the kernel reclaim its pages. A raw mapped relay must wait for send-side
reuse before returning its receive slots. On an earlier build, the mapped-object
relay used 15\% less relay-host CPU per delivered byte than \splice at a fixed
1\GBs (three runs per path). Figure~\ref{fig:tcp} reserves separate panels for
the receive, send and relay comparisons on the final build.
\todom{Populate each TCP panel by size and worker count. Match request shape,
show achieved rate and uncertainty, and identify fallback copies.}
A separate kernel-bypass baseline, F-Stack~\cite{fstack}, runs TCP in user
space. Comparing the same relay on F-Stack requires counting polling cores,
reserved buffer memory and any copies needed to retain the incoming bytes.
\todom{Add the F-Stack relay with its actual buffer ownership, byte verification,
NIC and polling policy. Report both resource use and delivered rate.}

\begin{figure*}[t]
\centering
% Panel of the final-build TCP comparison. Populate from the qualified campaign.
\begin{minipage}[t]{0.315\textwidth}
\centering
\figplaceholder{2.20cm}{\textbf{(a) Capture}\\
Copy, named and mapped objects\\
Raw slots and discard controls\\[4pt]
\emph{Results pending}}
\end{minipage}
\hfill
% Panel of the final-build TCP comparison. Populate from the qualified campaign.
\begin{minipage}[t]{0.315\textwidth}
\centering
\figplaceholder{2.20cm}{\textbf{(b) Send}\\
Copy, per-send pin, registered pool\\
Named and mapped objects\\[4pt]
\emph{Results pending}}
\end{minipage}
\hfill
% Panel of the final-build TCP comparison. Populate from the qualified campaign.
\begin{minipage}[t]{0.315\textwidth}
\centering
\figplaceholder{2.20cm}{\textbf{(c) Forward}\\
Copy, splice, named and mapped objects\\
Raw mapped relay\\[4pt]
\emph{Results pending}}
\end{minipage}

\caption{TCP experiments. Each panel relates host CPU per delivered GB and
goodput to the range size. Capture and forward use received data. Send separates
prepared data from objects captured from a stream.}
\label{fig:tcp}
\end{figure*}

\paragraph{Device support.} On the tested virtio-net interfaces, Linux hardware
\recvzc registration returned \code{EOPNOTSUPP}. Named capture ran with the
ordinary driver, and mapped capture ran with the modified provider. This test
establishes availability on those interfaces. A comparison with hardware
\recvzc requires a NIC that supports its queue and header-placement operations.
\todom{Add the supported-NIC comparison and tested driver configurations;
distinguish integration availability from measured cost.}

\subsection{UDP: cost per delivered datagram}
\label{sec:eval:udp}

A datagram relay must preserve message boundaries and select a destination for
each message. Its useful output is the unique payload received at the sink.
Counting bytes accepted by a send would reward a relay that drops traffic
downstream. The comparison uses the same IPv4 relay policy, offered load,
datagram sizes and queue bounds for copying \iouring, named objects and mapped
objects. For prefix routing, each relay examines the same bytes. An in-place
edit is a separate workload with a copying relay that makes the same edit.

Figure~\ref{fig:udp} pairs CPU per delivered GB with datagrams delivered per
second and loss at each offered load. Small datagrams expose per-message work
that can disappear in a bulk byte rate. A request/response case measures the
latency of that decision at a common delivered rate.
\begin{figure}[t]
\centering
\figplaceholder{2.25cm}{UDP results pending\\
Prefix routing and same-length edits\\
CPU per delivered GB, datagrams/s and loss}
\caption{UDP relay comparison. Counting complete, unique datagrams at the sink
makes loss visible even when the sender maintains its offered rate.}
\label{fig:udp}
\end{figure}
\todom{Fill the UDP figure and report the observed effect for each mode by size
and load, with delivered rates, uncertainty and request/response latency.}

\subsection{Block I/O: serve and ingest}
\label{sec:eval:block}

\paragraph{Read and serve.} A storage server reads aligned chunks from a raw
device window and sends them over TCP. The baseline uses direct reads into
registered application buffers followed by \sendzc. The object paths read
into kernel-owned buffers and send the returned handles. Sending each chunk
once measures a transfer. Serving the same chunk to several clients also
exercises retention. All paths use the same chunk size, physical read policy,
queue depth and number of recipients. Useful bytes are counted at the clients.
A file-cache hit followed by \sendfile is a distinct workload because it can
avoid the device read entirely.

\paragraph{Receive and write.} In the reverse direction the baseline receives
into a registered buffer and issues a direct write from it. An object write
accepts only a compatible layout. The dense-capture prototype copies incoming
TCP data once into compatible pages. That copy belongs in the cost of ingest.
For mapped access the comparison requires writable pages that both endpoints
accept. Here the denominator is bytes in completed device writes. A separate
readback checks their contents. Neither configuration requests durable writes.

\begin{figure}[t]
\centering
\figplaceholder{2.25cm}{Block I/O results pending\\
Device read to socket send\\Socket receive to device write}
\caption{Storage transfers in both directions. Read-and-serve counts payload
delivered to clients. Ingest counts completed write bytes. Each comparison
pairs server CPU cost with useful throughput and device bytes transferred.}
\label{fig:block}
\end{figure}
Figure~\ref{fig:block} separates the two directions so a saved copy on the read
path cannot obscure a required copy on the write path.
\todom{Insert named and mapped results for both directions, including the
registered baseline, copies, physical I/O, depth and uncertainty. Identify
the exact device and whether reuse changes the number of physical reads.}

\subsection{Application costs and memory pressure}
\label{sec:eval:rr}
\label{sec:eval:store}
\label{sec:eval:retention}
\label{sec:eval:planned}

Bulk transfers amortize object creation over many bytes. A request/response
client instead waits for each reply, making per-object work visible. A direct
connection provides a latency reference. For each relay the comparison varies
message size and concurrent clients separately and reports median and tail
latency, transactions per second and relay CPU per transaction.
\todom{Insert TCP request/response results at common loads, including small
messages and one- and multi-worker configurations.}

A key-value server holds data across requests. Its copying, named and mapped
backends use the same parser, keys and eviction policy. The relevant capacity is
useful value bytes under the same \emph{physical} memory budget. A small slice
can hold a large page, and a slow client can keep an evicted value's pages busy.
Reporting live payload, backing-page slack, metadata and reserved pool space
explains where the memory goes. Shared physical pages count once. Throughput
and tail latency under that budget show the price of retaining more values.
\todom{Insert capacity, goodput, latency and reclamation results for all three
backends, including fan-out and slow readers. Include CPU-local recycling
costs and reserved capacity at matched worker counts.}

For mapped receive, releasing a handle can leave a page waiting for TCP
retransmission. A slow-reader test can distinguish pages waiting on such
references from pages ready to return to the device. An independent protection
test changes header bytes in the mapping after the kernel takes its private
copy and checks which bytes the stack uses. \reviewold{These tests exercise the two
promises of segregation: the mapper may edit data, and the kernel retains
control of parsing and page reuse.\space}%
\reviewnew{A planned lifetime test queues a packet on another service's socket,
lets the mapper exit, then checks the payload at the receiver.}
\todom{Add the private-header and buffer-reclamation checks, including
exhaustion, cross-CPU returns and teardown with outstanding references.}
